% !TeX root = ../main.tex
\section{Appendix: Parameter Provenance}
\label{sec:appendix-provenance}

This appendix consolidates parameter provenance and the targeted sensitivity checks used to bound the main claims. Whenever a value is not directly reported in the source literature, we label it explicitly as a proxy estimate or analytical assumption.

\subsection{Three-Seed Summary for the Canonical V4 Path}
\label{subsec:v4-rerun-sanity}

To document the reproducibility of the canonical Tiny-ViT HAT path, we lock the three-seed V4 summary directly in the appendix. Each seed was evaluated with a \(10\)-run noisy Monte Carlo sweep under the canonical deployment regime (4-bit quantization with 5\% C2C and 10\% D2D variability), and the final cross-seed statistic is computed from the three seed means.

\begin{table}[ht]
\centering
\caption{Three-seed reproducibility summary for the canonical V4 regime.}
\label{tab:v4-three-seed-summary}
\begin{tabular}{lcc}
\toprule
\textbf{Seed / Aggregate} & \textbf{Noisy MC Accuracy} & \textbf{Statistic Type} \\
\midrule
42 & \(87.64\pm0.48\)\% & 10-run mean \(\pm\) std \\
123 & \(88.10\pm0.33\)\% & 10-run mean \(\pm\) std \\
2026 & \(88.11\pm0.47\)\% & 10-run mean \(\pm\) std \\
\midrule
Cross-seed aggregate & \(\mathbf{87.95\pm0.27}\)\% & mean \(\pm\) std of seed means \\
\bottomrule
\end{tabular}
\end{table}

\FloatBarrier

\begin{table}[ht]
\centering
\caption{Parameter provenance tracking matrix.}
\label{tab:provenance}
\scriptsize
\setlength{\tabcolsep}{1.8pt}
\renewcommand{\arraystretch}{1.1}
\begin{tabular}{>{\raggedright\arraybackslash}p{2.1cm}>{\raggedright\arraybackslash}p{1.9cm}>{\raggedright\arraybackslash}p{2.2cm}>{\raggedright\arraybackslash}p{2.0cm}>{\raggedright\arraybackslash}p{5.0cm}}
\toprule
\textbf{Parameter} & \textbf{Canonical Organic Profile} & \textbf{Zhang 2025 OPECT Case Study} & \textbf{Task 34-36 Stress Tests} & \textbf{Provenance / Derivation Notes} \\
\midrule
 Conductance Range (\(G_{\max}/G_{\min}\)) & 10 & 47.3 & 10 & Canonical anchored to scalable OPECTs; Zhang derives from reported max/min current levels \citep{zhang2025opect}. \\ 
 Effective States ($n_{states}$) & 16 (4-bit) & 34 & 16 (4-bit) & Canonical matches general low-precision target. Zhang anchored to Fig.3h \& Supp.Fig.8 \citep{zhang2025opect}. \\ 
 Cycle-to-Cycle Noise (\(\sigma_{\text{C2C}}\)) & 5\% & 2\% & 5\% (Proportional) & Zhang value is a proxy estimate from 8-cycle repeatability (Supp.Fig.15) \citep{zhang2025opect}. Proportional stress scales \(\sigma \propto |G|\). \\
  Device-to-Device Variability (\(\sigma_{\text{D2D}}\)) & 10\% & 3\% & 10\% (Proportional) & Zhang uses 3\% as a conservative conductance-domain proxy from the reported $\sim$1\% $V_{th}$ spread \citep{zhang2025opect}. \\
 Retention Decay (\(\tau_1, \tau_2, A_0\)) & 140ms, 610ms, 0.6 & Not fitted & 140ms, 610ms, 0.6 & Canonical dual-exponential fit anchored to Vincze \citep{vincze2025dualplasticity}. Zhang (Fig.2d) is qualitative only. \\ 
 Nonlinear Write Asymmetry (\(NL\)) & 1.0 (LTP), -1.0 (LTD) & Canonical defaults retained & 2.0 (LTP), -2.0 (LTD) & Zhang does not report a pulse-level NL fit, so the case study keeps the canonical symmetric-write defaults rather than injecting a separate literature-fit NL term. Stress tests (\(NL=2\)) enforce severe gradient-scaling asymmetry. \\ 
 Noise Mode & Uniform & Uniform & Proportional & Uniform injects noise normalized to full range. Proportional applies state-dependent noise magnitudes. \\ 
\bottomrule
\end{tabular}
\end{table}

\FloatBarrier

\subsection{Vincze et al. 2025 Parameter Extraction}
\label{subsec:vincze-extraction}

Because Vincze et al.\ \citep{vincze2025dualplasticity} is currently available as a 2025 Early Access paper, we explicitly list the extracted parameters used to define our canonical organic profile. These values were derived from the reported transient and steady-state characterization of the dual-plasticity organic synaptic transistor:

\begin{itemize}
    \item \textbf{Conductance Window}: $G_{\max}/G_{\min} \approx 10$, anchored to the scalable operation regime.
    \item \textbf{Retention Time Constants}: The dual-exponential fit for temporal drift yields a fast decay constant $\tau_1 = 140$~ms, a slow decay constant $\tau_2 = 610$~ms, and an initial amplitude ratio $A_0 = 0.6$. These align with the short-term and long-term plasticity dynamics reported for the device.
    \item \textbf{Nonlinear Write Asymmetry ($NL$)}: While the canonical profile uses symmetric writes ($NL=1.0$), the physical profile derived from the literature indicates moderate write nonlinearity corresponding to $NL_{LTP} = 2.0$ and $NL_{LTD} = -2.0$, which we evaluate directly in our severe stress tests.
\end{itemize}

By explicitly defining these parameters in the configuration files, the simulation framework allows any reader to reproduce our exact device assumptions without requiring access to the original materials manuscript.

\subsection{Proxy Estimate Sensitivity Analysis}
\label{subsec:proxy-sensitivity}

The Zhang 2025 OPECT values for \(\sigma_{\text{C2C}}\) and \(\sigma_{\text{D2D}}\) are transparent proxy estimates rather than directly measured full-array noise distributions. We therefore ran a compact 2D sensitivity sweep on the Ensemble HAT checkpoint to verify that the case-study conclusion does not hinge on one narrow proxy choice.

\begin{table}[ht]
\centering
\caption{Ensemble HAT accuracy under varied Zhang proxy noise assumptions. The repeated rows across the C2C sweep are expected under the present checkpoint because the deployment accuracy is dominated by static D2D variability; the sampled C2C term contributes below the reported precision after Monte Carlo averaging.}
\label{tab:sensitivity}
\small
\begin{tabular}{lccccc}
\toprule
\textbf{C2C \textbackslash D2D} & \textbf{2\%} & \textbf{3\% (Nominal)} & \textbf{5\%} & \textbf{10\%} & \textbf{15\%} \\
\midrule
\textbf{1\%} & 88.57\% & 88.53\% & 88.33\% & 87.30\% & 84.59\% \\
\textbf{2\% (Nominal)} & 88.57\% & 88.53\% & 88.33\% & 87.30\% & 84.59\% \\
\textbf{5\%} & 88.57\% & 88.53\% & 88.33\% & 87.30\% & 84.59\% \\
\textbf{8\%} & 88.57\% & 88.53\% & 88.33\% & 87.30\% & 84.59\% \\
\bottomrule
\end{tabular}
\end{table}

\begin{table}[ht]
\centering
\caption{Statistical summary of the C2C-insensitivity trend in Table~\ref{tab:sensitivity}. The confidence interval uses a normal approximation over the 10 Monte Carlo runs recorded for the nominal C2C \(=2\%\) row.}
\label{tab:sensitivity-ci}
\begin{tabular}{lccc}
\toprule
\textbf{D2D} & \textbf{Nominal Acc.} & \textbf{95\% CI} & \(\mathbf{\Delta_{\max}}\) \textbf{across C2C sweep} \\
\midrule
2\%  & 88.57\% & [88.52, 88.62]\% & 0.00 pp \\
3\%  & 88.53\% & [88.48, 88.58]\% & 0.00 pp \\
5\%  & 88.33\% & [88.27, 88.38]\% & 0.00 pp \\
10\% & 87.30\% & [87.22, 87.38]\% & 0.00 pp \\
15\% & 84.59\% & [84.51, 84.66]\% & 0.00 pp \\
\bottomrule
\end{tabular}
\end{table}

\textit{Interpretation.} Table~\ref{tab:sensitivity-ci} makes the same point more formally: at fixed D2D mismatch, the maximum accuracy excursion across the full C2C sweep is 0.00 percentage points at the reported precision, while the Monte Carlo uncertainty band remains narrow. The numerically identical rows in Table~\ref{tab:sensitivity} should therefore be read as a stability result rather than as a copy-paste artifact: for this checkpoint and evaluation budget, static D2D spatial variability dominates the case-study outcome, whereas per-forward C2C sampling noise remains sub-resolution. Even under a pessimistic D2D assumption of 15\%, the model remains above 80\% accuracy, far from the 10\% collapse of the standard V4 model.

\FloatBarrier

\subsection{Uniform vs. State-Dependent Retention Comparison}
\label{subsec:retention-comparison}

To test whether the first-order uniform retention assumption is sufficient for the present paper, we compared it with a more physical state-dependent model in which high-conductance states decay up to 20\% faster than low-conductance states. The comparison uses the V4 Ensemble HAT model on CIFAR-10, because that checkpoint is the retained deployment-facing path emphasized in the main text.

\begin{table}[ht]
\centering
\caption{Retention accuracy comparison: uniform vs. state-dependent decay.}
\label{tab:retention-comparison}
\begin{tabular}{lccc}
\toprule
\textbf{Time (s)} & \textbf{Uniform (\%)} & \textbf{State-Dep (\%)} & \textbf{Diff (pp)} \\
\midrule
0    & 90.77 & 90.80 & +0.03 \\
1    & 90.07 & 90.08 & +0.01 \\
10   & 89.78 & 89.81 & +0.03 \\
100  & 89.87 & 89.89 & +0.02 \\
1000 & 89.82 & 89.90 & +0.08 \\
\bottomrule
\end{tabular}
\end{table}

The accuracy difference stays below 0.1 percentage points across all tested time intervals. For this paper, the result serves only as a regime-specific sanity check, not as a claim that state dependence is negligible in general. We did not replicate the same comparison on the collapsing fresh-instance standard-HAT path or on non-HAT checkpoints, so this table should be read as a scoped validation of the \emph{retention-model choice} within the retained Ensemble-HAT deployment regime rather than as a universal statement across all training recipes.

\FloatBarrier

\subsection{Automated Profile Fitting Pipeline}
\label{subsec:auto-fitting}

To reduce manual profile tuning in future measured-device studies, we implemented an automated fitting utility (\path{profile_auto_fitter_gpt.py}). It ingests four raw-data categories:
\begin{enumerate}
    \item \textbf{Pulsed Programming Curves}: Fits the conductance range ($G_{\min}, G_{\max}$) and nonlinearity coefficients ($NL_{LTP}, NL_{LTD}$) using a behavioral exponential model.
    \item \textbf{Cycle-to-Cycle (C2C) Statistics}: Extracts the $\sigma_{\text{C2C}}$ noise parameter from repeated write-verify measurements at identical target states.
    \item \textbf{Device-to-Device (D2D) Mismatch}: Calculates the $\sigma_{\text{D2D}}$ parameter from conductance distributions across multiple physical devices in an array.
    \item \textbf{Retention Decay}: Fits the double-exponential decay constants ($\tau_1, \tau_2, A_0$) from long-term conductance-vs-time measurements.
\end{enumerate}

The output is a standardized JSON device profile that can be injected directly into the PyTorch inference and HAT loops. In the current paper, this utility is included as supporting infrastructure for follow-on calibration workflows rather than as a core source of the reported headline results. A detailed calibration protocol is provided in Supplementary Note~S-HW.
